% Source: physical PDF pages 73--77; printed pages 50--54.
% Appendix A begins after the final Section 15.9 footnote on physical p. 73
% and ends before the Appendix B heading on physical p. 77.
% Figures: none. Source uncertainties: none.
\chapterappendix{A}{A Theorem Regarding Lie Algebras}

In this appendix we consider a general Lie algebra $\mathcal G$, with generators $t_\alpha$ and structure constants $C^\alpha{}_{\beta\gamma}$, and prove the equivalence of three conditions:

\begin{description}
\item[\textbf{a:}] There exists a real symmetric positive-definite matrix $g_{\alpha\beta}$ that satisfies the invariance condition
\begin{equation}
 g_{\alpha\beta}C^\beta{}_{\gamma\delta}
 =-g_{\gamma\beta}C^\beta{}_{\alpha\delta}.
 \tag{15.A.1}\label{eq:15.A.1}
\end{equation}
(This is the condition \eqref{eq:15.2.4} shown in Section 15.2 to be necessary on physical grounds.)

\item[\textbf{b:}] There is a basis for the Lie algebra (that is, a set of generators $\widetilde t_\alpha=\mathcal S_{\alpha\beta}t_\beta$, with $\mathcal S$ a real non-singular matrix) for which the structure constants $\widetilde C^\alpha{}_{\beta\gamma}$ are antisymmetric not only in the lower indices $\beta$ and $\gamma$ but in all three indices $\alpha$, $\beta$ and $\gamma$.

\item[\textbf{c:}] The Lie algebra $\mathcal G$ is the direct sum of commuting compact simple and $U(1)$ subalgebras $\mathcal H_m$.
\end{description}

We shall prove the equivalence of statements \textbf{a}, \textbf{b}, and \textbf{c} by showing that \textbf{a} implies \textbf{b}, \textbf{b} implies \textbf{c}, and \textbf{c} implies \textbf{a}. As a by-product, we shall also show that if these conditions are satisfied, then it is possible to choose the generators of $\mathcal G$ as $t_{ma}$, with $m$ labelling the simple or $U(1)$ subalgebra $\mathcal H_m$ to which $t_{ma}$ belongs and $a$ labelling the individual generators in this subalgebra, and with the matrix $g_{ma,nb}$ that satisfies Eq.~\eqref{eq:15.A.1} taking the form
\begin{equation}
 g_{ma,nb}=g_m^{-2}\delta_{mn}\delta_{ab},
 \tag{15.A.2}\label{eq:15.A.2}
\end{equation}
where the $g_m^{-2}$ are arbitrary real positive constants.

First, let us assume \textbf{a}, the existence of a real symmetric positive-definite $g_{\alpha\beta}$ satisfying the invariance condition \eqref{eq:15.A.1}. We may then define new generators
\begin{equation}
 \widetilde t_\alpha\equiv(g^{-1/2})_{\alpha\beta}t_\beta,
 \tag{15.A.3}\label{eq:15.A.3}
\end{equation}
in which the existence of a real inverse square-root matrix $g^{-1/2}$ is guaranteed by the positive-definiteness of $g_{\alpha\beta}$. These satisfy a Lie algebra
\begin{equation}
 [\widetilde t_\alpha,\widetilde t_\beta]
 =i\widetilde C_{\alpha\beta\gamma}\widetilde t_\gamma,
 \tag{15.A.4}\label{eq:15.A.4}
\end{equation}
where
\begin{equation}
 \widetilde C_{\gamma\alpha\beta}
 \equiv(g^{-1/2})_{\alpha\alpha'}(g^{-1/2})_{\beta\beta'}
 (g^{+1/2})_{\gamma\gamma'}C^{\gamma'}{}_{\alpha'\beta'}.
 \tag{15.A.5}\label{eq:15.A.5}
\end{equation}
(In this basis it is convenient to drop the distinction between upper and lower indices $\alpha$, $\beta$, etc., and write $\widetilde C_{\gamma\alpha\beta}$ in place of $\widetilde C^\gamma{}_{\alpha\beta}$.) Then Eq.~\eqref{eq:15.A.1} tells us that $\widetilde C_{\alpha\gamma\delta}$ is antisymmetric in $\alpha$ and $\gamma$ as well as in $\gamma$ and $\delta$, and hence is totally antisymmetric, verifying \textbf{b}.

Next let us assume \textbf{b}, the existence of a basis for the Lie algebra for which the structure constants are totally antisymmetric. In this basis, the matrices $(\widetilde t^A_\alpha)_{\beta\gamma}=-i\widetilde C_{\beta\gamma\alpha}$ of the adjoint representation are imaginary and antisymmetric, and hence Hermitian. According to a general theorem about Hermitian matrices,\footnote{If a set of matrices $H_\alpha$ is not irreducible, then by definition there must be a set of vectors $u_n$ that span a subspace (other than the whole space) that is left invariant by the $H_\alpha$: that is, for all $\alpha$ and $n$, $H_\alpha u_n=\sum_m(C_\alpha)_{mn}u_m$. In this case we can adopt a basis consisting of the vectors $u_n$ together with a set $v_k$ that span the space orthogonal to all the $u_n$. If the $H_\alpha$ are Hermitian then $(u_m,H_\alpha v_k)=\sum_n(C_\alpha)_{mn}(u_n,v_k)=0$, so the space spanned by the $v_k$ is also invariant under the $H_\alpha$: $H_\alpha v_k=\sum_l(D_\alpha)_{lk}v_l$. In this basis the matrices $H_\alpha$ are simultaneously reduced to a block-diagonal form:
\[
 H_\alpha=\begin{pmatrix}C_\alpha&0\\0&D_\alpha\end{pmatrix}.
\]
Continuing in this way, we can completely reduce the matrices $H_\alpha$ to a block-diagonal form with irreducible matrices in the blocks.} the $\widetilde t^A_\alpha$ are then either irreducible or totally reducible.

By an irreducible set of $N\cdot N$ matrices $\widetilde t^A_\alpha$ is meant a set for which there exists no subspace of dimensionality $<N$ that is left invariant by all the $\widetilde t^A_\alpha$---that is, no set of less than $N$ non-zero vectors $(u_r)_\beta$ for which $(\widetilde t^A_\alpha)_{\beta\gamma}(u_r)_\gamma$ is for each $\alpha$ and $r$ a linear combination of vectors $(u_s)_\beta$. Since the matrices $(\widetilde t^A_\alpha)_{\beta\gamma}$ are proportional to the structure constants in this basis, this is equivalent to the statement that there is no set of linear combinations $\mathcal T_r\equiv(u_r)_\gamma\widetilde t^A_\gamma$ that is closed under commutation with all the $\widetilde t^A_\alpha$, that is, for which $[\widetilde t^A_\alpha,\mathcal T_r]$ is for each $\alpha$ and $r$ a linear combination of the $\mathcal T_s$. Such a set of matrices $\mathcal T_r$ would furnish the generators of an invariant subalgebra of the full Lie algebra; the absence of such a set means that the Lie algebra is simple.

By a totally reducible set of matrices $\widetilde t^A_\alpha$ is meant one that by a suitable choice of basis may be written as block-diagonal supermatrices
\begin{equation}
 (\widetilde t^A_\alpha)_{ma,nb}
 =[\widetilde t^{A(m)}_\alpha]_{ab}\delta_{mn},
 \tag{15.A.6}\label{eq:15.A.6}
\end{equation}
where the submatrices $\widetilde t^{A(m)}_\alpha$ are either irreducible or vanish.\footnote{Here $m$ and $n$ label the blocks along the main diagonal, and $a$ and $b$ label rows and columns within these blocks. Also, $m$ is not summed, and the range of the indices $a$, $b$, etc. in general depends on $m$.} Adopting this basis also for the Lie algebra itself, the structure constants are then
\begin{equation}
 \widetilde C_{\ell c,ma,nb}
 =i(\widetilde t^A_{\ell c})_{ma,nb}
 =i(\widetilde t^{A(m)}_{\ell c})_{ab}\delta_{mn}.
 \tag{15.A.7}\label{eq:15.A.7}
\end{equation}
But since this is totally antisymmetric, and proportional to $\delta_{mn}$, it must also be proportional to $\delta_{\ell n}$ and $\delta_{\ell m}$:
\begin{equation}
 \widetilde C_{\ell c,ma,nb}
 =\delta_{\ell n}\delta_{\ell m}C^{(\ell)}_{cab}.
 \tag{15.A.8}\label{eq:15.A.8}
\end{equation}
In other words, for any representation $t_a^{(m)}\equiv t_{ma}$ of the Lie algebra in this basis, we have
\begin{equation}
 [t_a^{(m)},t_b^{(n)}]
 =i\delta_{mn}C^{(m)}_{cab}t_c^{(m)},
 \tag{15.A.9}\label{eq:15.A.9}
\end{equation}
with $C^{(m)}_{cab}$ real and totally antisymmetric in the indices $a$, $b$, $c$. The fact that it is possible to construct a basis in which the generators fall into sets $t^{(m)}$, with the commutators of the generators in one set with each other given by a linear combination of generators in that set, and with all members of one set commuting with all members of any other set, is what we mean when we say that the Lie algebra is a \emph{direct sum} of the subalgebras $t^{(m)}$. For each $m$, the set of matrices of the adjoint representation $t^{(m)A}$ is either irreducible or zero, corresponding to a subalgebra that is either simple or consists of so-called $U(1)$ generators, which commute with all generators of the whole algebra.

We have thus shown that the most general Lie algebra with totally antisymmetric real structure constants is a direct sum of one or more simple and/or $U(1)$ Lie algebras. Furthermore, the simple subalgebras are compact, in the sense that each matrix $-C^{(m)}_{acd}C^{(m)}_{bdc}$ is positive-definite, because for any real vector $u_a$, $-C^{(m)}_{acd}C^{(m)}_{bdc}u_au_b=\sum_{cd}[\sum_a u_aC^{(m)}_{acd}]^2$ is a sum of positive quantities, which cannot vanish unless $u_a=0$, since for $u_a\ne0$ the condition that $\sum_a u_aC^{(m)}_{acd}=0$ would imply that $\sum_a u_at_a^{(m)}$ is itself an invariant Abelian subalgebra, in contradiction with the fact that the $t_a^{(m)}$ form a simple Lie algebra. This completes the proof of \textbf{c}.

Finally, let us assume \textbf{c}, that we have a Lie algebra that is the direct sum of a set of simple or $U(1)$ Lie algebras; that is, in some basis $t_a^{(m)}=\mathcal S_{ma,\alpha}t_\alpha$ with real non-singular $\mathcal S$, we have
\[
 [t_a^{(m)},t_b^{(n)}]
 =i\delta_{mn}C^{(m)c}{}_{ab}t_c^{(m)},
\]
where each subalgebra $t^{(m)}$ is either simple or commutes with everything. We furthermore assume that the simple subalgebras are compact, in the sense that the matrices
\begin{equation}
 g_{ab}^{(m)}\equiv-C^{(m)c}{}_{ad}C^{(m)d}{}_{bc}
 \tag{15.A.10}\label{eq:15.A.10}
\end{equation}
are positive-definite. To construct a real positive-definite matrix $g_{\alpha\beta}$ satisfying Eq.~\eqref{eq:15.A.1}, in the basis of generators $t_{ma}=t_a^{(m)}$, we take
\begin{equation}
 g_{ma,nb}=g_{ab}^{(m)}\delta_{mn},
 \tag{15.A.11}\label{eq:15.A.11}
\end{equation}
where $g_{ab}^{(m)}$ is taken as the matrix \eqref{eq:15.A.10} when $t^{(m)}$ is a simple subalgebra, and is taken as an arbitrary real symmetric positive-definite matrix when $t^{(m)}$ is a direct sum of one or more $U(1)$ subalgebras. The matrix \eqref{eq:15.A.11} is obviously real, symmetric, and positive-definite because each $g_{ab}^{(m)}$ is. To check the requirement \eqref{eq:15.A.1}, recall the Jacobi identity for the structure constants:
\begin{equation}
 C^{(m)c}{}_{ad}C^{(m)d}{}_{be}
 +C^{(m)c}{}_{bd}C^{(m)d}{}_{ea}
 +C^{(m)c}{}_{ed}C^{(m)d}{}_{ab}=0
 \tag{15.A.12}\label{eq:15.A.12}
\end{equation}
and contract with $C^{(m)e}{}_{fc}$. After renaming the summation indices in the third term so that $c\longrightarrow d\longrightarrow e\longrightarrow c$ and using Eq.~\eqref{eq:15.A.10}, the result for the simple subalgebras may be written:
\[
 g_{df}^{(m)}C^{(m)d}{}_{ab}
 =C^{(m)c}{}_{ad}C^{(m)d}{}_{be}C^{(m)e}{}_{cf}
 -C^{(m)c}{}_{fd}C^{(m)d}{}_{be}C^{(m)e}{}_{ca}.
\]
The important point is that this shows that the left-hand side is antisymmetric in $a$ and $f$:
\begin{equation}
 g_{df}^{(m)}C^{(m)d}{}_{ab}
 =-g_{da}^{(m)}C^{(m)d}{}_{fb}.
 \tag{15.A.13}\label{eq:15.A.13}
\end{equation}
The same result holds trivially for the $U(1)$ subalgebras, where the structure constants vanish. The symmetry condition \eqref{eq:15.A.1} follows immediately from Eq.~\eqref{eq:15.A.13}, thus completing the proof of \textbf{a}. This concludes our proof of the equivalence of statements \textbf{a}, \textbf{b}, and \textbf{c}.

Now we come back to the matrix $g_{\alpha\beta}$. With totally antisymmetric structure constants, the invariance condition \eqref{eq:15.2.4} may be expressed as the statement that this matrix commutes with all the matrices in the adjoint representation of the Lie algebra
\begin{equation}
 [g,\widetilde t^A_\gamma]=0.
 \tag{15.A.14}\label{eq:15.A.14}
\end{equation}
We have seen that all $(\widetilde t^A_\gamma)_{\alpha\beta}$ can be put in block-diagonal form, with irreducible (or zero) submatrices along the main diagonal. A well-known theorem~\hyperref[ch15-ref-25]{[25]} tells us that then $g_{\alpha\beta}$ must also be block-diagonal, with blocks of the same size and position as in the $\widetilde t^A_\gamma$, and with the submatrix in each block proportional to the unit matrix. (Where two of the submatrices in the $\widetilde t^A_\gamma$ are equivalent, in the sense that they can be related by a similarity transformation, it may be necessary to make a suitable change of basis in order to bring the submatrices of $g_{\alpha\beta}$ into a form proportional to unit matrices.) In the notation of Eq.~\eqref{eq:15.A.11}, the metric is then given by Eq.~\eqref{eq:15.A.2}.
